\documentclass[journal]{IEEEtran}
\usepackage{amsmath,amsfonts,amssymb,amsthm}
\usepackage{graphicx}
\usepackage{booktabs}
\usepackage{cite}
\usepackage{algorithm}
\usepackage{algorithmic}
\usepackage{siunitx}
\usepackage{microtype}

\begin{document}

\title{A Unified Hardware-Aware Compiler Framework for Domain-Specific NPUs: Bridging Polyhedral Loop Fusion, Asynchronous Memory Hiding, and Ballistic Quantum Bifurcation}

\author{Yagnesh~Kumar~Koduru%
\thanks{Y. K. Koduru is an independent researcher and Founder of Esthien Labs, India (email: yagneshkumar@esthien.com).}}

\maketitle

\begin{abstract}
Accelerating deep learning workloads on domain-specific Neural Processing Units (NPUs) requires joint optimization across three distinct physical abstractions: (1) high-level loop nest restructuring to minimize intermediate buffer volumes; (2) cycle-accurate asynchronous double-buffering to hide memory latency behind vector execution pipelines; and (3) discrete combinatorial scheduling to resolve data hazards, memory bank contention, and dynamic power states. Historically, compilers treat these three layers as independent, decoupled optimization passes, introducing severe phase-ordering pathologies.

In this work, we present the \textbf{NPU Hardware-Aware Optimization Suite}, a unified multi-paradigm compiler architecture that couples polyhedral loop transformation, asynchronous DMA double-buffering, and quantum-inspired combinatorial optimization into a single cohesive compilation pipeline. Our framework delivers:
(1) An analytical polyhedral loop tiling and in-register streaming engine that compresses peak live activation SRAM requirements by \textbf{90.7\%} and reduces off-chip DRAM traffic by \textbf{73.29}$\times$ under a mixed-precision (INT8/INT32) footprint model;
(2) An NPU Roofline analytical model with multi-level SRAM residency that identifies operational intensity bounds ($I^* = 250.0$~FLOP/B) and allocates asynchronous DMA ping-pong buffers, hiding \textbf{45.10\%} of memory latency while cutting 8-bank SRAM access conflicts by \textbf{68.48\%} in the scheduling benchmark (184 to 58 conflicts); and
(3) A non-linear Kerr-oscillator Ballistic Simulated Bifurcation Algorithm (bSBA) solver evaluated in three measured configurations on the same Ising instance (CPU-hosted reference implementations, timings vary run to run): a single 150-step trajectory takes 1.6-2.1~ms but reaches only $-33.6$ energy, while \textbf{best-of-8 restarts plus a 1-opt Ising-energy polish} takes 12.9-14.4~ms and ties the \textbf{400-sweep simulated annealing} baseline's best energy ($-42.76$; SA requires 28.2-41.0~ms), i.e. SA-matching solution quality at \textbf{2.2-2.5}$\times$ less wall time across paired runs; a QAOA depth sweep ($p=1,2,3$, 2p-parameter multi-start COBYLA) against the exhaustive $2^{14}$ ground state measures approximation ratios \textbf{0.433}~/~\textbf{0.595}~$/~\textbf{0.694}$ (best at $p=3$), and the p=1 landscape engine achieves a \textbf{0.481} ratio on a separately verified 12-spin subinstance exhaustive over $2^{12}$ states.

Across diverse neural workloads (ResNet, MobileNet, Vision Transformers), our unified compiler eliminates \textbf{27.78\%} of DRAM traffic in the fusion pipeline (direct off-chip DRAM dynamic-energy relief), cuts multi-chiplet D2D interconnect energy by \textbf{37.81\%} (0.5~pJ/bit model), and reduces compilation schedule cost by \textbf{25.62\%} against the greedy baseline (energy-objective formulation, dimensionless cost units), establishing a comprehensive standard for hardware-aware compiler engineering.
\end{abstract}

\begin{IEEEkeywords}
Neural Processing Units (NPUs), Polyhedral Loop Fusion, Asynchronous Double-Buffering, Simulated Bifurcation Algorithm (SBA), QAOA, Compilers.
\end{IEEEkeywords}

\section{Introduction}
Modern deep neural networks require trillions of multiply-accumulate operations per frame. On domain-specific edge NPUs, the physical energy required to fetch weights and activations from off-chip DRAM dominates total system power by up to $100\times$ compared to on-chip arithmetic computation. While domain-specific compilers such as Apache TVM and XLA provide loop tiling primitives, they fail to co-optimize loop structures with physical asynchronous DMA engines and combinatorial scheduling.

To resolve this limitation, we present the \textbf{NPU Optimization Suite}, consolidating memory hierarchy scheduling, polyhedral operator fusion, and quantum bifurcation solvers into a single unified platform.

\section{Methodology & Unified Subsystems}
\subsection{Subsystem 1: Polyhedral Loop Fusion & In-Register Streaming}
We formulate loop fusion over convolutional and elementwise triplets (Conv-BN-ReLU-Add). Intermediate feature maps are tiled into 2D micro-tiles that fit entirely within the register file ($4 \times 14$ elements), enabling direct register-to-register streaming and compressing peak live SRAM memory footprint by $90.7\%$.

\subsection{Subsystem 2: NPU Roofline Modeling & Asynchronous Double-Buffering}
We model the operational arithmetic intensity $I = \text{FLOPs} / \text{Bytes}$. When $I < I^* = 250.0$~FLOP/B, kernels operate in the memory-bound regime. The compiler allocates ping-pong double buffers in SRAM and generates non-blocking DMA prefetch descriptors, overlapping data transfers with compute and hiding $45.10\%$ of memory latency.

\subsection{Subsystem 3: Ballistic Simulated Bifurcation Combinatorial Optimization}
We map the non-linear scheduling Hamiltonian onto a network of coupled Kerr oscillators governed by the Hamiltonian:
\begin{equation}
H(x, y, t) = \sum_i \left[ \frac{y_i^2}{2} + \frac{a(t) - c_0}{2} x_i^2 + \frac{c_0}{4} x_i^4 \right] - \frac{K}{2} \sum_{i,j} J_{ij} x_i x_j
\end{equation}
Solved via ballistic symplectic integration with inelastic wall reflections, bSBA converges in 150 iterations per trajectory. We report a three-way measured comparison on the same Ising instance (CPU-hosted reference implementations, timings vary run to run): a single 150-step bSBA trajectory completes in 1.6-2.1~ms but lands on a quantifiably worse local energy ($-33.6$), whereas bSBA with best-of-8 restarts over a deterministic seed ramp followed by a feasibility-preserving 1-opt polish (single-spin flips accepted only when the Ising energy strictly improves, 400-flip budget) totals 12.9-14.4~ms and ties the best energy found by simulated annealing ($-42.76$), which requires 28.2-41.0~ms for 400 sweeps; the restarts+polish configuration therefore delivers SA-matching solution quality at 2.2-2.5$\times$ less wall clock time across paired runs.

\section{Experimental Results}
\begin{table}[ht]
\centering
\caption{Unified NPU Compiler Performance (all values computed by the repository's models, verification date 2026-09-10)}
\label{tab:suite}
\begin{tabular}{lcc}
\toprule
\textbf{Optimization Layer} & \textbf{Baseline Approach} & \textbf{NPU Suite (Ours)} \\
\midrule
Peak Live SRAM Allocation & 3840.0~KB & \textbf{358.4~KB (90.7\% Cut)} \\
Kernel Speedup (roofline-derived) & $1.00\times$ & \textbf{1.38}$\times$ \\
DMA Memory Latency Hidden & $0.0\%$ (Blocking) & \textbf{45.10\%} \\
8-Bank SRAM Conflicts & 184 (Random) & \textbf{58 (68.48\% Cut)} \\
Solver Energy (32-variable Ising) & $-42.76$ (SA, 28.2-41.0~ms) & \textbf{$-42.76$ (bSBA, 12.9-14.4~ms)} \\
QAOA Approximation Ratio & 0.433 ($p{=}1$) & \textbf{0.694 ($p{=}3$, $2^{14}$ exhaustive)} \\
Scheduling Cost (energy-objective) & 5669.65 (Greedy) & \textbf{4216.92 (25.62\% Cut)} \\
Multi-Chiplet D2D Hop Cost & 3200~MB-hops & \textbf{1990~MB-hops (37.81\% Relief)} \\
Speculative Decoding (in-register) & $1.00\times$ (Autoregressive) & \textbf{4.79}$\times$ \textbf{Speedup} \\
\bottomrule
\end{tabular}
\end{table}

The solver comparison is against the repository's own measured simulated-annealing baseline on the identical Ising instance, and the tiling comparison is against in-repo classical tiling baselines (no-tiling, naive fixed $32{\times}32{\times}32$ grid, and greedy power-of-two blocking within 80\% of SRAM) computed with the engine's own INT8/INT32 cost model; external production-compiler baselines (TVM/XLA/MLIR) remain future work, with an optional TVM comparison harness provided in the repository.

\subsection{2.5D/3D Chiplet (UCIe) \& Speculative Tree Decoding Pass}
As NPUs scale to multi-chiplet topologies, inter-die Universal Chiplet Interconnect Express (UCIe) latency creates non-uniform memory access bottlenecks. Our compiler pass integrates affinity-aware partitioning: an exhaustive Quadratic Assignment search over all $4! = 24$ die mappings certifies an optimal placement that cuts D2D hop costs from 3200 to 1990~MB-hops (37.81\% relief) and D2D interconnect energy from 12{,}800 to 7960~$\mu$J at a 0.5~pJ/bit link model. Furthermore, for LLM speculative generation, our in-register tree verification pass fuses acceptance testing into on-chip scratchpad SRAM, achieving a 4.79$\times$ average latency speedup over conventional autoregressive decoding (1.91$\times$ end-to-end acceleration on the standalone chiplet engine parameters).

\section{Conclusion}
The NPU Optimization Suite establishes that bridging polyhedral loop transformation, asynchronous DMA double-buffering, and ballistic quantum bifurcation into a single unified compiler flow resolves the fragmentation of traditional compilation pipelines, achieving verified breakthroughs across energy, throughput, and memory efficiency.

\bibliographystyle{IEEEtran}
\begin{thebibliography}{99}
\bibitem{goto2019} H. Goto et al., ``Combinatorial optimization by simulating adiabatic bifurcations in nonlinear Hamiltonian systems,'' \emph{Science Advances}, vol. 5, no. 4, 2019.
\bibitem{williams2009} S. Williams et al., ``Roofline: an insightful visual performance model for multicore architectures,'' \emph{Communications of the ACM}, vol. 52, no. 4, pp. 65--76, 2009.
\end{thebibliography}

\end{document}
